% !TEX root = ../../main.tex
%% Défi du chapitre 15 — L'intégrale de Gauss
%% Résultats vérifiés par verifications/defis_11_15.py.

\rubrique{Défi}

\begin{defi}{L'intégrale de Gauss}{2 h}
La fonction $t\mapsto\mathrm{e}^{-t^{2}}$ n'a pas de primitive exprimable à
l'aide des fonctions usuelles (théorème de Liouville, admis) ; on va
pourtant calculer exactement l'aire totale sous sa courbe, la plus célèbre
des probabilités. On pose
\[
F(A)=\int_0^{A}\mathrm{e}^{-t^{2}}\,\mathrm{d}t\quad(A\geqslant0),
\qquad
W_k=\int_0^{\pi/2}\cos^{k}\theta\,\mathrm{d}\theta\quad(k\in\N),
\]
ces dernières étant les intégrales de Wallis, dont on rappelle
l'équivalent $W_k\sim\sqrt{\pi/(2k)}$ (chapitre 11, exercice 23).

\partieExo{A}{Existence et encadrement}
\begin{questions}
  \item Montrer que $F$ est croissante et que $F(A)\leqslant1+\frac1{\mathrm{e}}$
    pour tout $A\geqslant0$. \emph{Comparer $\mathrm{e}^{-t^{2}}$ et
    $\mathrm{e}^{-t}$ sur $\intFO{1}{+\infty}$.} En déduire que $F(A)$ admet
    une limite finie $G$ quand $A\to+\infty$.
  \item Montrer que $1+u\leqslant\mathrm{e}^{u}$ pour tout réel $u$. En
    déduire, pour $n\geqslant1$ et $t\in\intFF{0}{\sqrt n}$ :
    \[
    \Bigl(1-\frac{t^{2}}{n}\Bigr)^{n}\leqslant\mathrm{e}^{-t^{2}}
    \leqslant\Bigl(1+\frac{t^{2}}{n}\Bigr)^{-n} .
    \]
\end{questions}

\medskip
\begin{center}
\begin{tikzpicture}[x=2.3cm,y=2.6cm,line cap=round]
  \fill[otGreen,opacity=0.18] (0,0) -- plot[domain=0:2.6,samples=90] ({\x},{exp(-\x*\x)}) -- (2.6,0) -- cycle;
  \draw[otMuted,line width=0.5pt,-{Stealth[length=1.6mm]}] (0,0) -- (2.8,0);
  \draw[otMuted,line width=0.5pt,-{Stealth[length=1.6mm]}] (0,0) -- (0,1.2);
  \draw[otNavy!60,line width=0.9pt,dash pattern=on 3pt off 1.5pt] plot[domain=0:2.6,samples=90]
    ({\x},{(1+\x*\x/3)^(-3)});
  \draw[otNavy,line width=0.9pt,dash pattern=on 1pt off 1.2pt] plot[domain=0:1.7320,samples=90]
    ({\x},{(1-\x*\x/3)^3});
  \draw[otGreenDark,line width=1.3pt] plot[domain=0:2.6,samples=90] ({\x},{exp(-\x*\x)});
  \draw[otMuted,line width=0.4pt] (1.7320,0.03) -- (1.7320,-0.03);
  \node[font=\tiny,text=otMuted,below] at (1.7320,-0.02) {$\sqrt3$};
  \node[font=\tiny,text=otMuted,anchor=east] at (0,1) {$1$};
  \node[font=\scriptsize,text=otGreenDark,anchor=south west] at (0.62,0.72) {$\mathrm{e}^{-t^{2}}$};
  \node[font=\scriptsize,text=otNavy!70,anchor=west] at (1.25,0.33) {$\bigl(1+\frac{t^{2}}{3}\bigr)^{-3}$};
  \node[font=\scriptsize,text=otNavy,anchor=east] at (0.5,0.28) {$\bigl(1-\frac{t^{2}}{3}\bigr)^{3}$};
  \draw[otNavy,line width=0.4pt] (0.5,0.28) -- (0.86,0.31);
\end{tikzpicture}\\[0.5mm]
\legende{La courbe de Gauss (en vert, aire $G$ ombrée) encadrée par les
deux fonctions de la question 2 pour $n=3$.}
\end{center}

\partieExo{B}{Le calcul par les intégrales de Wallis}
\begin{questions}[resume]
  \item Par le changement de variable $t=\sqrt n\sin\theta$, montrer que
    \[
    \int_0^{\sqrt n}\Bigl(1-\frac{t^{2}}{n}\Bigr)^{n}\mathrm{d}t
    =\sqrt n\,W_{2n+1} .
    \]
  \item Par le changement de variable $t=\sqrt n\tan\theta$, montrer que,
    pour $n\geqslant2$,
    \[
    \int_0^{\sqrt n}\Bigl(1+\frac{t^{2}}{n}\Bigr)^{-n}\mathrm{d}t
    \leqslant\sqrt n\,W_{2n-2} .
    \]
  \item En déduire $\sqrt n\,W_{2n+1}\leqslant F\bigl(\sqrt n\bigr)\leqslant\sqrt n\,W_{2n-2}$,
    puis
    \[
    G=\lim_{A\to+\infty}\int_0^{A}\mathrm{e}^{-t^{2}}\,\mathrm{d}t=\frac{\sqrt\pi}{2} .
    \]
\end{questions}

\partieExo{C}{Applications}
\begin{questions}[resume]
  \item Pour $a>0$, déterminer la limite de
    $\displaystyle\int_0^{A}\mathrm{e}^{-at^{2}}\,\mathrm{d}t$ quand
    $A\to+\infty$.
  \item Montrer que
    $\dfrac{1}{\sqrt{2\pi}}\displaystyle\int_{-A}^{A}\mathrm{e}^{-t^{2}/2}\,\mathrm{d}t
    \xrightarrow[A\to+\infty]{}1$ : c'est la normalisation de la \emph{loi
    normale}, limite universelle des sommes de variables aléatoires
    indépendantes (programme de deuxième année).
  \item D'où vient le $\sqrt\pi$ ? Comparer avec les apparitions de $\pi$
    dans la formule de Stirling (chapitre 15) et dans la marche aléatoire
    (chapitre 17).
\end{questions}
\end{defi}
